\Needspace{5\baselineskip}
\SourceLeaf{422}%
\begin{center}6.\end{center}
\noindent\hspace*{1.2em}Betrachtet man $x$ und $y$ als Funktionen von $s$, so nehmen die Formeln eine noch elegantere Gestalt an. Man setze
\[
\begin{aligned}
\frac{dx}{ds}&=x',&\qquad \frac{dx'}{ds}&=x'',\\
\frac{dy}{ds}&=y',&\frac{dy'}{ds}&=y''.
\end{aligned}
\]
Dann gilt
\[
\begin{aligned}
x'&=\cos\varphi,&\qquad y'&=\sin\varphi,\\
x''&=-\frac{\sin\varphi}{r},&y''&=\frac{\cos\varphi}{r},
\end{aligned}
\]
oder
\[
y'=-rx'',\qquad x'=ry''
\]
oder auch
\[
1=r(x'y''-y'x''),
\]
sodass
\[
x'y''-y'x''
\]
die Krümmung darstellt und
\[
\frac{1}{x'y''-y'x''}
\]
den Krümmungsradius.

\Needspace{5\baselineskip}
\begin{center}7.\end{center}
\noindent\hspace*{1.2em}Wir gehen jetzt zur Betrachtung der krummen Flächen über. Um die Richtungen gerader Linien im dreidimensionalen Raum zu veranschaulichen, denken wir uns um einen beliebigen Mittelpunkt eine Kugeloberfläche mit dem Radius $1$. Ein Punkt auf ihr stellt somit die Richtung aller geraden Linien dar, die parallel zu dem Radius sind, an dessen Endpunkt er liegt. So wie die Lagen aller Punkte im Raum durch ihre senkrechten Abstände $x,y,z$ von drei aufeinander senkrechten Ebenen bestimmt werden, so sollen die Richtungen der drei Hauptachsen, die auf diesen Hauptebenen senkrecht stehen, auf der Oberfläche der Hilfskugel\space%
